\scanpage{96}
\mspage{47}
donnée $G$, l'ens $\widetilde A(G)$ et $I(G)$ (munis l'un et autre d'une application \inserted{``origine''} à valeurs dans $S(G)$).

Pour ceci, on \unsure{a|à|va} utiliser la notion de ``graphes avec branches aplaties'', qui généralis\unsure{e} celle de graphe, et topologiquement correspond à la classification des \unsure{\deleted{4}|4|\deleted{q}|q} couples $(X,S)$ d'un espace topologique compact $X$, et d'une partie finie $S$, tels que a) $X$ soit une étoile \unsure{au|en} voisinage de $\forall$ pt de $S$, \deleted{\unsure{et|a}} b) \deleted{\unsure{S}} tt composante connexe de $X\setminus S$ est une $1$-variété à bord non compacte (car on veut que tout \deleted{\unsure{adhérence} d'} ``arête'' ait un sommet). \unsure{Il|Ça} revient au $=$ de dire que $X$ est un $1$-complexe top, et $S$ une partie finie qui contient les sommets d'ordre $\geq3$, et telle qu'aucune composante connexe de $X\setminus S$ ne soit compacte (i.e. que tt composante connexe de $X$ contienne un sommet, i.e. $S\unsure{\rightharpoondown|{-}}\unsure{\Pi|\pi}_0(X)$ est surjectif).

Ceci se classifie combinatoirement par des systèmes $(S,\unsure{\widehat A|\widehat I},\sigma_{\unsure{\widehat A|\widehat I}},o_{\widetilde A})$, où \deleted{$\unsure{\widetilde A=\sigma_{\widetilde A}|\widetilde A,\sigma_A|\widetilde A,o_A}$} $S,\unsure{\widehat A|\widehat I}$ sont des ens finis, $\sigma_{\unsure{\widetilde A|\widehat A}}$ est une involution de $\widetilde A$, et $o_{\widetilde A}:\widetilde A\unsure{\rightharpoondown|{-}}S$ une
\margin{95}
